% Section 3: Normalization
% Focus: Training stability, optimization landscape, and the evolving understanding of what normalization truly does.

Normalization is arguably the most inconspicuous yet consequential component of the Transformer architecture.
Unlike the attention mechanism, which defines the model's identity, or position encoding, which shapes its inductive biases, normalization operates as a silent guardian of training stability---its presence is rarely questioned, but its absence is immediately catastrophic.
The choice of normalization type, placement, and scope directly determines the maximum trainable depth, hardware efficiency, and scaling behavior of large language models.
From the introduction of Batch Normalization~\cite{ioffe2015batch} in 2015 to the recent emergence of normalization-free architectures such as Dynamic Tanh~\cite{zhu2025transformers} and Dynamic Erf~\cite{chen2025stronger}, the evolution of normalization reflects an ongoing quest to understand what training stability fundamentally requires.
This section traces nine technical threads that collectively map this evolution: from foundational methods to efficient variants, from placement theory to attention-internal constraints, from conditional modulation to holistic architectures, and from automated search to the radical proposal that normalization layers may be dispensable altogether.

%% 3.1 Foundation: BN -> LN -> IN -> GN
\subsection{Foundation: From Batch Normalization to Layer Normalization}

The foundational normalization methods established the principle that normalizing along a chosen set of dimensions stabilizes training by controlling the distribution of activations.
Batch Normalization (BN)~\cite{ioffe2015batch} pioneered this paradigm by computing mean and variance across the batch dimension: given an activation $x_i$, BN computes $\hat{x}_i = (x_i - \mu_\mathcal{B}) / \sqrt{\sigma^2_\mathcal{B} + \epsilon}$ followed by a learnable affine transformation $y_i = \gamma \hat{x}_i + \beta$.
This design enabled higher learning rates, reduced initialization sensitivity, and provided implicit regularization, fueling the success of deep convolutional networks.
However, BN's reliance on batch-level statistics renders it fundamentally incompatible with Transformer-based language models, where variable sequence lengths cause batch statistics to fluctuate dramatically.

Layer Normalization (LN)~\cite{ba2016layer} resolved this incompatibility by normalizing across the feature dimension $d$ within each individual sample, making the computation entirely independent of batch composition.
Given a $d$-dimensional hidden vector $\mathbf{h}$, LN computes:
\begin{equation}
    \text{LN}(\mathbf{h}) = \gamma \odot \frac{\mathbf{h} - \mu}{\sqrt{\sigma^2 + \epsilon}} + \beta, \quad \mu = \frac{1}{d}\sum_{i=1}^{d} h_i, \quad \sigma^2 = \frac{1}{d}\sum_{i=1}^{d}(h_i - \mu)^2,
\end{equation}
where $\gamma, \beta \in \mathbb{R}^d$ are learnable scale and shift parameters.
This formulation provides three decisive advantages for language modeling: (i)~complete independence from batch size, enabling single-sample inference identical to training; (ii)~no need for running statistics, eliminating the train--test discrepancy inherent in BN; and (iii)~compatibility with distributed training without cross-device synchronization.
Instance Normalization~\cite{ulyanov2016instance}, which normalizes per-channel spatial statistics, found its niche in style transfer but has negligible relevance to NLP.
Group Normalization (GN)~\cite{wu2018group} partitions channels into $G$ groups and normalizes within each group, providing a continuous interpolation between Instance Normalization ($G=C$) and Layer Normalization ($G=1$).
GN has recently resurfaced in the Differential Transformer~\cite{ye2024differential}, where it normalizes per-head outputs after the differential attention subtraction---a role that leverages GN's group-wise independence.
The migration from BN to LN established feature-dimension normalization as the default paradigm for sequence models, setting the baseline against which all subsequent innovations are measured.

%% 3.2 Efficient Variants: RMSNorm, ScaleNorm, PowerNorm
\subsection{Efficient Normalization Variants}

Having established Layer Normalization as the standard, a natural question arises: which components of LN are truly necessary?
RMSNorm~\cite{zhang2019root} provides a definitive partial answer by hypothesizing that the re-scaling invariance of LN, rather than its re-centering (mean subtraction), is the key property responsible for training stability.
RMSNorm simplifies LN by eliminating the mean computation entirely:
\begin{equation}
    \text{RMSNorm}(\mathbf{a}) = \frac{\mathbf{a}}{\text{RMS}(\mathbf{a})} \odot \gamma, \quad \text{RMS}(\mathbf{a}) = \sqrt{\frac{1}{n}\sum_{i=1}^{n} a_i^2},
\end{equation}
where $\gamma \in \mathbb{R}^n$ is a learnable gain (no bias $\beta$).
This eliminates one reduction operation and halves the number of learnable parameters from $2d$ to $d$, yielding measured speedups of 7--64\% depending on implementation.
The hypothesis has been substantiated by subsequent geometric analysis~\cite{gupta2024reintroducing}, which reveals that LN's mean subtraction is equivalent to projecting out the component along the all-ones vector $\mathbf{1} = [1, \ldots, 1]^\top$---an irreversible operation that discards information.
Crucially, LLMs trained with RMSNorm (such as LLaMA) naturally learn representations orthogonal to $\mathbf{1}$, rendering the explicit mean subtraction redundant.
RMSNorm is now the de facto standard in virtually all frontier LLMs, including LLaMA~1/2/3, PaLM, Gemma, Qwen, Mistral, and DeepSeek.

ScaleNorm~\cite{nguyen2019transformers} pursues even more extreme simplification: $\text{ScaleNorm}(\mathbf{x}) = g \cdot \mathbf{x} / \|\mathbf{x}\|_2$, using a single global scalar $g$ in place of per-element gains.
This maps all vectors to a hypersphere of radius $g$, foreshadowing the hyperspherical geometry later developed by nGPT~\cite{loshchilov2024ngpt}.
PowerNorm~\cite{shen2020powernorm} takes a different tack, diagnosing why BN fails in NLP---batch statistics fluctuate 3--5$\times$ more in text than in vision---and replacing per-batch variance with an exponential moving average of the quadratic mean.
Despite its theoretical insights, PowerNorm's residual dependence on the batch dimension prevented mainstream adoption.
The trajectory from LN to RMSNorm exemplifies a minimality principle: normalization should retain only the invariance properties that matter, discarding unnecessary computations.

%% 3.3 Normalization Placement: Pre-Norm, Post-Norm, and Beyond
\subsection{Normalization Placement: Pre-Norm, Post-Norm, and Beyond}

Where normalization is placed relative to the sublayer is among the most consequential architectural decisions in Transformer design.
The original Transformer~\cite{vaswani2017attention} adopted \textbf{Post-LN}, applying normalization after the residual addition: $\mathbf{x}_{l+1} = \text{LN}(\mathbf{x}_l + \text{Sublayer}(\mathbf{x}_l))$.
Mean field theory analysis~\cite{xiong2020layer} revealed that Post-LN suffers from large gradient norms near the output layers at initialization, necessitating careful learning rate warm-up.
\textbf{Pre-LN} reverses the order---$\mathbf{x}_{l+1} = \mathbf{x}_l + \text{Sublayer}(\text{LN}(\mathbf{x}_l))$---creating a ``gradient highway'' through which gradients flow unattenuated across arbitrary depth.
Pre-LN eliminates the need for warm-up and has become the default in most modern LLMs, though some evidence suggests Post-LN achieves slightly better final performance when training is successful.

Several proposals seek to combine the stability of Pre-LN with the performance of Post-LN.
\textbf{Sandwich-LN}~\cite{ding2021cogview} applies normalization both before and after the sublayer, originally developed to suppress NaN losses during FP16 training of the 4B-parameter CogView model; Gemma~2~\cite{gemma2024improving} later adopted this pattern with Pre+Post RMSNorm.
\textbf{DeepNorm}~\cite{wang2024deepnet} scales the residual connection by a constant $\alpha > 1$: $\mathbf{x}_{l+1} = \text{LN}(\alpha \cdot \mathbf{x}_l + f(\mathbf{x}_l))$, coupled with theoretically derived initialization scaling.
\begin{equation}
x_{l+1} = \mathrm{LN}\!\left(\alpha \cdot x_l + F_l(x_l)\right), \quad \alpha = (2L)^{1/4},\; \text{init: } W \sim \mathcal{N}(0, \beta^2), \; \beta = (8L)^{-1/4}
\label{eq:deepnorm}
\end{equation}
This enabled training a 1,000-layer Transformer (2,500 sublayers), demonstrating that a 200-layer 3.2B model can outperform a 48-layer 12B model by 5 BLEU points---evidence that depth is an efficient scaling axis.
\textbf{Sub-LN / MAGNETO}~\cite{wang2022magneto} pushes normalization inside sublayers at a finer granularity, inserting additional LN between projections within both attention and FFN blocks, achieving uniform improvements across language, vision, speech, and multimodal tasks.
\textbf{Peri-LN}~\cite{kim2025periln} applies normalization at both ends of the sublayer while keeping the residual path direct: $\mathbf{x}_{l+1} = \mathbf{x}_l + \text{LN}_2(f(\text{LN}_1(\mathbf{x}_l)))$, achieving controlled linear variance growth $\text{Var}(\mathbf{x}_{l+1}) \approx \text{Var}(\mathbf{x}_l) + O(1)$.

More recent work pushes beyond fixed placement entirely.
\textbf{GeoNorm}~\cite{li2026geonorm} replaces traditional normalization with geodesic updates on a representation manifold, $\mathbf{h}_{l+1} = \text{Exp}_{\mathbf{h}_l}(\alpha_l \cdot \mathbf{v}_l)$, unifying Pre-LN and Post-LN as special cases of tangent-space transformations and manifold projections, and introducing a normalization schedule analogous to learning rate schedules.
\textbf{SiameseNorm}~\cite{wang2025siamesenorm} adopts a dual-stream architecture maintaining both a Pre-Norm stream (for gradient stability) and a Post-Norm stream (for representation quality) with shared parameters, fusing their outputs at inference---though at the cost of nearly doubled computation.
A unifying perspective comes from particle interaction theory~\cite{karagodin2025normalization}, which models token representations as interacting particles on a hypersphere: self-attention acts as an inter-particle force, while normalization acts as a velocity regulator, with different placements corresponding to different velocity--stability trade-offs.

%% 3.4 Attention-Internal Normalization
\subsection{Attention-Internal Normalization}

While traditional normalization operates at the sublayer level, the attention mechanism's internal computations harbor their own sources of numerical instability---most critically, the unbounded growth of $QK^\top$ dot-product logits as model depth and hidden dimension increase.
\textbf{QK-Norm}~\cite{henry2020query} addresses this by $\ell_2$-normalizing the query and key vectors before the dot product:
\begin{equation}
    \hat{Q} = Q / \|Q\|_2, \quad \hat{K} = K / \|K\|_2, \quad \text{Attn} = \text{softmax}\!\left(g \cdot \hat{Q}\hat{K}^\top\right) V,
\end{equation}
where $g$ is a per-head learnable scalar initialized to $\sqrt{d_k}$.
By constraining logits to cosine similarities in $[-1, 1]$ and then rescaling by $g$, QK-Norm fundamentally eliminates logit explosion and prevents softmax saturation.
This technique has been adopted by Gemma~2/3, Qwen~3, OLMo~2, and Chameleon, though it is notably incompatible with DeepSeek-V3's Multi-Latent Attention, which does not materialize full Q and K tensors.

\textbf{Logit soft-capping}~\cite{gemma2024improving} provides a complementary output-side constraint: $\text{logits} \leftarrow c \cdot \tanh(\text{logits} / c)$, bounding logits within $(-c, +c)$.
Gemma~2 employs attention logit capping at $c = 50$ and final logit capping at $c = 30$, creating a multi-layered defense.
In the Differential Transformer~\cite{ye2024differential}, the subtraction of two softmax maps produces outputs whose values are no longer bounded by $[0, 1]$, necessitating GroupNorm~\cite{wu2018group} applied per attention head to restore controlled activation scales; the $(1 - \lambda_{\text{init}})$ rescaling factor further aligns gradient flow.
Recent work on Lp-QKNorm~\cite{lopez2026enhanced} generalizes L2 normalization to arbitrary $L_p$ norms, demonstrating that different values of $p$ induce different attention space geometries, though validation remains limited to small-scale settings.
Attention-internal normalization has transitioned from an optional enhancement to a necessary safeguard at scale, with Gemma~2's triple-layered protection (Pre+Post RMSNorm, QK-Norm, soft-capping) representing the current state of engineering practice.

%% 3.5 Conditional and Adaptive Normalization
\subsection{Conditional and Adaptive Normalization}

Standard normalization treats the affine parameters $\gamma$ and $\beta$ as static, shared across all inputs.
Conditional normalization makes these parameters functions of an external conditioning signal, transforming the normalization layer into a lightweight but powerful mechanism for signal injection.
This idea traces back to FiLM (Feature-wise Linear Modulation)~\cite{perez2018film}, which modulates visual features with question embeddings in visual question answering via $\text{FiLM}(\mathbf{x}) = \gamma(\mathbf{c}) \odot \mathbf{x} + \beta(\mathbf{c})$.
AdaNorm~\cite{xu2019understanding} explored input-dependent normalization strength within Transformers, adapting the degree of normalization based on the ``difficulty'' of the input, though it did not achieve widespread adoption.

The decisive breakthrough came with \textbf{adaLN-Zero}~\cite{peebles2023scalable} in the Diffusion Transformer (DiT).
Here, conditioning signals (timestep embedding and class embedding) are passed through a shared MLP to dynamically generate six parameter sets per block---scale $\gamma$, shift $\beta$, and gate $\alpha$ for both the attention and FFN sublayers:
\begin{equation}
    (\gamma_1, \beta_1, \alpha_1, \gamma_2, \beta_2, \alpha_2) = \text{MLP}\!\left(\text{embed}(t) + \text{embed}(c)\right),
\end{equation}
\begin{equation}
    \text{output} = \mathbf{x} + \alpha \odot \text{Sublayer}\!\left(\gamma \odot \frac{\mathbf{x} - \mu}{\sigma} + \beta\right).
\end{equation}
The zero-initialization of the MLP's output layer ensures that each block initially acts as the identity function, providing training stability analogous to Fixup~\cite{zhang2019fixup}.
adaLN-Zero achieved state-of-the-art FID of 2.27 on ImageNet 256$\times$256 and significantly outperformed both cross-attention and in-context conditioning.
Its influence is pervasive: Stable Diffusion~3, Flux, and Sora all build upon the DiT architecture with adaLN conditioning.
This success establishes normalization layers as the canonical interface for conditioning signal injection in generative models, a pattern that may extend to cross-modal fusion in multimodal architectures.

%% 3.6 Holistic Normalization Architectures
\subsection{Holistic Normalization Architectures}

While the preceding methods treat normalization as a localized operation at specific positions, a more radical approach elevates normalization to a global architectural principle.
\textbf{nGPT}~\cite{loshchilov2024ngpt} constrains all vectors in the network---embeddings, hidden states, query/key/value projections, MLP weights, and output weights---to lie on the unit hypersphere via $\ell_2$ normalization.
Layer updates become spherical interpolations:
\begin{equation}
    \mathbf{h}_{l+1} = \text{Normalize}\!\left(\mathbf{h}_l + \alpha_l \cdot \left(f_l(\mathbf{h}_l) - \mathbf{h}_l\right)\right),
\end{equation}
where $\alpha_l$ is a learnable step size and the final normalization reprojects onto the sphere.
Equivalently, constraining all vectors to the unit hypersphere and using spherical linear interpolation:
\begin{equation}
\hat{x}_l = \frac{x_l}{\|x_l\|_2}, \quad x_{l+1} = \mathrm{slerp}\!\left(\hat{x}_l,\, \hat{F}_l(\hat{x}_l),\, \alpha_l\right)
\label{eq:ngpt}
\end{equation}
Because both Q and K are normalized, attention logits are naturally cosine similarities in $[-1, 1]$, eliminating the need for separate QK-Norm.
Critically, nGPT contains no traditional LayerNorm or RMSNorm layers---normalization is an intrinsic property of the representation space rather than an extrinsic correction.
The key geometric insight is that in unconstrained Euclidean space, a substantial fraction of optimizer updates merely change vector norms without altering directions, wasting computation.
The hyperspherical constraint eliminates this redundant degree of freedom, and nGPT achieves 4--20$\times$ reduction in steps to reach equivalent loss, though large-scale validation beyond nanoGPT remains pending.

\textbf{Gemma~2}~\cite{gemma2024improving} pursues the opposite philosophy---maximal redundancy for maximal robustness.
It deploys normalization at multiple levels simultaneously: Pre+Post RMSNorm at the sublayer level, QK-Norm at the attention input level, and logit soft-capping at both the attention and final output levels.
This multi-layered defense eliminates loss spikes across 2B, 9B, and 27B parameter scales.
\textbf{UnitNorm}~\cite{zhao2024unitnorm}, developed independently for time-series Transformers, applies pure $\ell_2$ normalization $\mathbf{x}/\|\mathbf{x}\|_2$ with zero additional parameters, addressing token shift and attention sparsity problems specific to temporal data.
Its convergence with nGPT's hyperspherical philosophy from entirely different problem domains suggests that projection onto the unit sphere may be a broadly effective inductive bias.
nGPT and Gemma~2 thus represent two polar philosophies of normalization design---geometric minimalism versus engineering defense-in-depth---and future architectures may need to navigate between these extremes.

%% 3.7 Automated Normalization-Activation Search
\subsection{Automated Normalization-Activation Search}

Normalization and activation functions are typically adjacent operations in Transformer blocks, and their functional roles---value range control and nonlinear transformation---exhibit significant overlap.
\textbf{EvoNorm}~\cite{liu2020evolving} challenges the assumption that these should be separately designed by defining a search space of primitive operations (addition, multiplication, division, sigmoid, variance, mean) and using evolutionary search to discover unified normalization-activation layers.
The search yielded two notable architectures: EvoNorm-B0 (batch-dependent), $\text{EvoNorm-B0}(\mathbf{x}) = \mathbf{x} \cdot \sigma(\mathbf{x} / v(\mathbf{x}))$, where sigmoid provides gating and variance provides normalization; and EvoNorm-S0 (batch-independent), which replaces variance with group standard deviation.
Both surpass their hand-designed counterparts (BN+ReLU and GN+ReLU, respectively) on ImageNet and COCO, and the layers discovered on small proxy tasks transfer directly to larger benchmarks.

The primary limitation of EvoNorm is its exclusive validation on convolutional architectures.
Whether the search paradigm can discover superior alternatives to the RMSNorm+SwiGLU combination now standard in Transformer FFN blocks remains an open question.
The search space would need to incorporate Transformer-specific primitives (e.g., attention logits, residual paths) and scale to architectures where evaluation cost is substantially higher.
Nevertheless, EvoNorm demonstrates that the joint normalization-activation design space is far from fully explored, and that automated methods can discover non-obvious functional compositions that outperform human intuition.

%% 3.8 Normalization-Free Architectures
\subsection{Normalization-Free Architectures}

If the essential functions of normalization can be achieved through other means---careful initialization, gradient clipping, or bounded activation functions---then normalization layers themselves may be unnecessary.
This line of inquiry began with \textbf{Fixup}~\cite{zhang2019fixup}, which showed that zero-initializing the last layer of each residual branch (making each block initially an identity function) controls variance growth and enables BatchNorm-free training of deep ResNets.
\textbf{T-Fixup}~\cite{huang2020tfixup} extended this to Transformers by scaling $W_V$ and $W_O$ by $(9N)^{-1/4}$ and initializing FFN output weights near zero, completely removing all LayerNorm layers and warm-up schedules while matching standard Transformer performance on machine translation.
\textbf{NFNet}~\cite{brock2021high} decomposed BN's implicit functions into two explicit mechanisms: Scaled Weight Standardization for controlling forward-pass signal variance, and Adaptive Gradient Clipping (AGC)---$G \leftarrow \lambda \frac{\|W\|_F}{\|G\|_F} G$ when $\|G\|_F / \|W\|_F > \lambda$---for constraining backward-pass update magnitude.
This achieved 86.5\% top-1 on ImageNet without any normalization.

The most disruptive recent development is \textbf{Dynamic Tanh (DyT)}~\cite{zhu2025transformers}, which observed that the input--output mapping of trained LayerNorm layers empirically resembles an S-shaped curve and proposed a direct replacement:
\begin{equation}
    \mathrm{DyT}(x) = \gamma \odot \tanh(\alpha \, x) + \beta, \quad \alpha \in \mathbb{R}^d \text{ (learnable per-dimension)}
    \label{eq:dyt}
\end{equation}
where $\alpha$ is a per-layer learnable scalar and $\gamma, \beta$ are per-channel parameters.
As a purely element-wise operation with no reduction, DyT is maximally hardware-friendly and matches normalization performance across ViT, LLaMA, DiT, and wav2vec~2.0.
\textbf{Dynamic Erf (Derf)}~\cite{chen2025stronger} refines this insight by analyzing the critical role of saturation speed: erf saturates more slowly than tanh (Gaussian decay $O(e^{-x^2})$ versus exponential decay $O(e^{-2|x|})$), preserving gradients at larger magnitudes---at $|x|=2$, the erf derivative is approximately 26$\times$ larger than tanh's.
Derf uses $\text{Derf}(\mathbf{x}) = \text{erf}(\alpha \mathbf{x} + s)$ and consistently outperforms both DyT and standard normalization across all tested architectures.
These results challenge a decade of normalization research by suggesting that the essential function of normalization may not be statistical standardization at all, but rather bounded nonlinear compression---a function that any sufficiently well-behaved sigmoidal mapping can provide.

%% 3.9 Frontier Directions
\subsection{Frontier Directions: Hardware Co-Design, Optimizer Coupling, and Outlier Dynamics}

The most recent research reveals that normalization is not an isolated architectural component but is deeply entangled with hardware precision, optimizer design, and emergent activation dynamics.

\paragraph{FP8 training and normalization.}
The transition to FP8 arithmetic exposes a fundamental tension: normalization layers that stabilize training in full precision can become error amplifiers in low precision~\cite{zhang2025whylowprecision}.
The mechanism involves a positive feedback loop: low-rank similarity structures emerge in attention, concentrating information in few dimensions; quantization errors in these critical dimensions are then amplified by RMSNorm's division (an underestimated RMS causes global over-scaling) and softmax's exponentiation, eventually triggering loss spikes.
FOG~\cite{sun2025fog} addresses this architecturally by mandating QK-Norm (bounding attention logits within FP8's dynamic range) and inserting additional RMSNorm at strategic positions to control FFN intermediate activation scales, achieving all-FP8 GEMM computation with 40\% throughput improvement over BF16.
The activation kurtosis serves as a practical monitoring metric: lower kurtosis indicates more FP8-friendly distributions.

\paragraph{Optimizer--normalization coupling.}
The Muon optimizer~\cite{jordan2025muon} reveals an unexpected duality: by orthogonalizing momentum updates via SVD decomposition ($\theta_t = \theta_{t-1} - \text{lr} \cdot UV^\top$), Muon implicitly maintains the spectral condition of weight matrices, controlling activation norms through the backward pass rather than the forward pass.
This is dual to normalization layers, which control forward-pass activation scales.
The combination of Muon (backward spectral control) with DyT/Derf (forward bounded compression) raises the intriguing possibility of completely eliminating traditional normalization from both passes---though this remains an open research direction.

\paragraph{Outlier--normalization symbiosis.}
Recent analysis~\cite{bondarenko2026outlier} reveals that the massive activations and attention sinks observed in trained Transformers are not pathological artifacts but functional features co-evolved with normalization.
When a single activation $x_k \gg x_i$ dominates, RMSNorm effectively becomes $\hat{x}_i \approx x_i \cdot \sqrt{d} / |x_k| \cdot \gamma_i$---the outlier's magnitude controls the effective scale of all other dimensions.
Removing normalization eliminates outliers but degrades performance; clipping outliers while retaining normalization is equally harmful.
This symbiosis has practical implications for quantization, which must preserve outlier functionality rather than suppress it, and provides a new interpretation of DyT/Derf as alternative rescaling mechanisms that do not depend on outlier co-evolution.

\paragraph{Module-level hybrid normalization.}
\textbf{HybridNorm}~\cite{chen2025hybridnorm} abandons the assumption that a single normalization strategy should apply uniformly, instead assigning QKV-Norm to attention modules (controlling logit range and output scale) and Post-Norm to FFN modules (preserving gradient signal quality).
This module-specific strategy proves effective for both dense Transformers and sparse Mixture-of-Experts architectures, where Post-Norm helps maintain consistent output scales across different experts.

\vspace{0.5em}
These frontier directions converge on a deeper insight: the essence of normalization may not reside in any specific network layer, but in maintaining certain invariants during training---spectral conditioning, norm boundedness, or value-range compression.
These invariants can be enforced through normalization layers (the traditional approach), bounded activation functions (DyT/Derf), optimizer design (Muon), architectural constraints (nGPT), or combinations thereof.

%% Comparison Table
\begin{table}[t]
\centering
\caption{Comparison of normalization methods for Transformer-based models. ``Norm Dim.''\ indicates the dimension along which statistics are computed. ``Params'' denotes learnable parameters per normalized dimension $d$. ``Reduction Ops'' counts the number of statistical reduction operations (mean, variance, or RMS) per forward pass.}
\label{tab:norm_comparison}
\small
\begin{tabular}{lccccl}
\toprule
\textbf{Method} & \textbf{Norm Dim.} & \textbf{Params} & \textbf{Reduction Ops} & \textbf{Year} & \textbf{Adoption} \\
\midrule
BatchNorm~\cite{ioffe2015batch} & Batch & $2d$ & 2 (mean, var) & 2015 & CNNs only \\
LayerNorm~\cite{ba2016layer} & Feature & $2d$ & 2 (mean, var) & 2016 & BERT, GPT-2 \\
RMSNorm~\cite{zhang2019root} & Feature & $d$ & 1 (RMS) & 2019 & LLaMA, Qwen, Mistral \\
ScaleNorm~\cite{nguyen2019transformers} & Feature & 1 & 1 (L2 norm) & 2019 & Limited \\
QK-Norm~\cite{henry2020query} & Per-head Q,K & 1/head & 1 (L2 norm) & 2020 & Gemma 2/3, Qwen 3 \\
adaLN-Zero~\cite{peebles2023scalable} & Feature & $6d$/block & 2 (mean, var) & 2023 & DiT, SD3, Sora \\
nGPT~\cite{loshchilov2024ngpt} & All vectors & $\alpha$/layer & 1 (L2 norm) & 2024 & Experimental \\
DyT~\cite{zhu2025transformers} & Element-wise & $2d{+}1$ & 0 & 2025 & Experimental \\
Derf~\cite{chen2025stronger} & Element-wise & 2 & 0 & 2025 & Experimental \\
\bottomrule
\end{tabular}
\end{table}
